% C058：先熟悉结构，再映射到抛物线。
\clearpage
\begin{explanationbox}
\textbf{第一步：只看一条直线上的相反方向。}
\begin{center}\begin{tikzpicture}[font=\small,>=stealth]
\draw[gray!45,rounded corners=3pt] (0,0) rectangle (14.2,12);
\node at (7.1,11.35) {先不画抛物线：同一直线，反向走，投影比互为相反数};

\begin{scope}[shift={(4.2,5.4)},scale=1.05]
\coordinate (P) at (0,0);\coordinate (C) at (3,4);\coordinate (D) at (-.75,-1);
\fill[blue!6] (P)--(C)--(3,0)--cycle;
\fill[orange!10] (P)--(D)--(-.75,0)--cycle;
\draw[gray] (-1.6,0)--(3.7,0);
\draw[blue!75!black,very thick] (P)--(C);
\draw[orange!85!black,very thick] (P)--(D);
\draw[dashed,gray] (C)--(3,0) (D)--(-.75,0);
\draw (2.8,0)--(2.8,.2)--(3,.2);
\draw (-.55,0)--(-.55,-.2)--(-.75,-.2);
\foreach \q in {P,C,D}\fill (\q) circle(2pt);
\node[above] at (C) {$C$};\node[below left] at (D) {$D$};\node[below right] at (P) {$P$};
\node[blue!75!black,above left] at (1.5,2) {$PC=5$};
\draw[<->,blue!75!black] (0,-2)--(3,-2) node[midway,fill=white] {$+3$};
\node[orange!85!black,left] at (-.8,-2) {$-3/4$};
\draw[<-,orange!85!black] (-.75,-1.6)--(0,-1.6);
\end{scope}
\node[anchor=west,align=left] at (9.4,8.4) {$PD=5/4$\\[12pt]向右投影为正\\[8pt]向左投影为负};
\node at (7.1,2.3) {$\displaystyle\frac{+3}{5}=-\frac{-3/4}{5/4}=\frac35$};
\node at (7.1,1.25) {长度始终为正，负号只出现在水平位移中。};
\node at (7.1,.55) {下一步对应：$P\to F,\ C\to A,\ D\to B$；横向比例记作 $t$。};

\end{tikzpicture}
\end{center}

先不记新公式：从中间点出发，两个方向的水平位移比互为相反数。
设 $t=(x_C-x_P)/PC$，则另一侧为 $-t$；一般非水平方向下 $-1<t<1$。
\end{explanationbox}
\clearpage
\begin{explanationbox}
\textbf{第二步：把中间点换成焦点，用准线给出两条焦半径。}
\begin{center}\begin{tikzpicture}[font=\small,>=stealth]
\draw[gray!45,rounded corners=3pt] (0,0) rectangle (14.2,12);
\node at (7.1,11.35) {放回抛物线：两端都用定义，把斜边换成到准线的距离};

\begin{scope}[shift={(3.6,4.9)},scale=.95]
\draw[teal!75!black,thick,samples=90,domain=-2.3:4.5,variable=\v] plot ({\v*\v/4},{\v});
\draw[gray,dashed,thick] (-1,-2.2)--(-1,5.1);
\draw[gray] (-1.5,0)--(5.5,0);

\coordinate (A) at (4,4);\coordinate (B) at (.25,-1);\coordinate (F) at (1,0);
\draw[blue!75!black,very thick] (F)--(A);
\draw[orange!85!black,very thick] (F)--(B);
\foreach \q in {A,B,F}\fill (\q) circle(3pt);
\node[above left=5pt] at (A) {$A$};
\node[below left=8pt] at (B) {$B$};\node[below right=6pt] at (F) {$F$};

\coordinate (H) at (-1,4);\coordinate (K) at (-1,-1);
\draw[blue!75!black,thick] (A)--(H);
\draw[orange!85!black,thick] (B)--(K);
\draw[gray] (-1,3.8)--(-.8,3.8)--(-.8,4);
\draw[gray] (-1,-.8)--(-.8,-.8)--(-.8,-1);
\fill (H) circle(2.5pt) node[left] {$H$};\fill (K) circle(2.5pt) node[left] {$K$};
\node[above] at (-1,5.1) {$\ell$};
\draw[gray,dashed] (-1,-2.2)--(-1,-2.7) (F)--(1,-2.7);
\draw[<->,violet] (-1,-2.6)--(1,-2.6) node[midway,below] {$p$};
\end{scope}
\node[anchor=west,align=left] at (9.4,8.2) {$AF=AH$（蓝）\\[12pt]$BF=BK$（橙）};
\node[anchor=west,align=left] at (9.4,5.8) {$\ell:x=-p/2$\\[12pt]$F=(p/2,0)$};
\node at (7.1,1.65) {$r=x_A+p/2,\qquad s=x_B+p/2$};
\node at (7.1,.65) {本图 $p=2$：$AH=AF=5$，$BK=BF=5/4$；两端使用同一条准线。};

\end{tikzpicture}
\end{center}

对应上一页的 $P,C,D$，现在分别是 $F,A,B$。
令 $r=AF,s=BF$。由 $x_F=p/2$，结合两端的定义式，
\[
 x_A-x_F=r-p,\qquad x_B-x_F=s-p.
\]
再用反向投影比：
\[
 \frac{r-p}{r}=-\frac{s-p}{s}=t
 \quad\Longrightarrow\quad
 \frac1r=\frac{1-t}{p},\qquad\frac1s=\frac{1+t}{p}.
\]
\textbf{倒数一加，方向项 $t$ 就消失，留下 $2/p$。}
$t$ 也可写作射线 $FA$ 的方向余弦，但不必先背极坐标公式。
\end{explanationbox}
\clearpage
\begin{explanationbox}
\textbf{第三步：让两段变长、变短，观察倒数怎样互补。}
\begin{center}\begin{tikzpicture}[font=\small,>=stealth]
\draw[gray!45,rounded corners=3pt] (0,0) rectangle (14.2,14);
\node at (7.1,13.35) {看见定值：左边把弦展开读长度，右边读倒数};

\node at (3.8,12.1) {原长度：蓝 $r$、橙 $s$};
\node at (10,12.1) {倒数：蓝 $1/r$、橙 $1/s$};
\node at (3.8,11.3) {$r=5,\ s=5/4$};
\draw[blue!75!black,very thick] (3.8,10.5)--(5.949999999999999,10.5);
\draw[orange!85!black,very thick] (3.8,10.5)--(3.2624999999999997,10.5);
\foreach \x/\name in {5.949999999999999/A,3.2624999999999997/B,3.8/F} {\fill (\x,10.5) circle(2pt);\node[below=7pt] at (\x,10.5) {$\name$};}
\fill[blue!30] (8,10.3) rectangle (8.8,10.8);
\fill[orange!40] (8.8,10.3) rectangle (12,10.8);
\draw[gray] (8,10.3) rectangle (12,10.8);
\node at (10,9.5) {$\displaystyle\frac1r+\frac1s=1$};
\node at (3.8,7.7) {$r=s=2$};
\draw[blue!75!black,very thick] (3.8,6.9)--(4.66,6.9);
\draw[orange!85!black,very thick] (3.8,6.9)--(2.94,6.9);
\foreach \x/\name in {4.66/A,2.94/B,3.8/F} {\fill (\x,6.9) circle(2pt);\node[below=7pt] at (\x,6.9) {$\name$};}
\fill[blue!30] (8,6.7) rectangle (10.0,7.2);
\fill[orange!40] (10.0,6.7) rectangle (12,7.2);
\draw[gray] (8,6.7) rectangle (12,7.2);
\node at (10,5.9) {$\displaystyle\frac1r+\frac1s=1$};
\node at (3.8,4.1) {$r=5/4,\ s=5$};
\draw[blue!75!black,very thick] (3.8,3.3)--(4.3374999999999995,3.3);
\draw[orange!85!black,very thick] (3.8,3.3)--(1.65,3.3);
\foreach \x/\name in {4.3374999999999995/A,1.65/B,3.8/F} {\fill (\x,3.3) circle(2pt);\node[below=7pt] at (\x,3.3) {$\name$};}
\fill[blue!30] (8,3.0999999999999996) rectangle (11.2,3.5999999999999996);
\fill[orange!40] (11.2,3.0999999999999996) rectangle (12,3.5999999999999996);
\draw[gray] (8,3.0999999999999996) rectangle (12,3.5999999999999996);
\node at (10,2.3) {$\displaystyle\frac1r+\frac1s=1$};

\node at (7.1,.6) {三行均取 $p=2$；短段对应较大的倒数，两块倒数合起来总是 $1$。};

\end{tikzpicture}
\end{center}

同一个 $p$ 对应同一个倒数总量。\textbf{图中的线段是原长度，色条是倒数，二者不能混读。}
通径对应中间一行：$r=s=p$，立刻复原 $1/p+1/p=2/p$。
\end{explanationbox}
\clearpage
\begin{explanationbox}
\textbf{第四步：把最短情形和退化情形一起认清。}
\begin{center}\begin{tikzpicture}[font=\small,>=stealth]
\draw[gray!45,rounded corners=3pt] (0,0) rectangle (14.2,10);
\node at (7.1,9.35) {同一条抛物线、同尺度：通径最短，水平线不是两端点弦};
\node at (2.35,8.6) {倾斜焦点弦};
\begin{scope}[shift={(2.35,5.3)},scale=.43]
\draw[teal!75!black,thick,domain=-2.5:4.2,samples=80,variable=\v] plot ({\v*\v/4},{\v});
\draw[->,gray] (-.7,0)--(4.7,0);
\fill (1,0) circle(5pt);\node at (2.1,-.7) {$F$};
\draw[blue!75!black,very thick] (1,0)--(4,4);\draw[orange!85!black,very thick] (1,0)--(.25,-1);\fill (4,4) circle(5pt);\fill (.25,-1) circle(5pt);
\end{scope}
\node[align=center] at (2.35,2.8) {$r=5,\ s=5/4$\\[12pt]$|AB|=25/4>4$};
\node at (7.1,8.6) {通径：可用且取等};
\begin{scope}[shift={(7.1,5.3)},scale=.43]
\draw[teal!75!black,thick,domain=-2.5:4.2,samples=80,variable=\v] plot ({\v*\v/4},{\v});
\draw[->,gray] (-.7,0)--(4.7,0);
\fill (1,0) circle(5pt);\node at (2.1,-.7) {$F$};
\draw[blue!75!black,very thick] (1,0)--(1,2);\draw[orange!85!black,very thick] (1,0)--(1,-2);\fill (1,2) circle(5pt);\fill (1,-2) circle(5pt);
\end{scope}
\node[align=center] at (7.1,2.8) {$r=s=p=2$\\[12pt]$|AB|=2p=4$};
\node at (11.85,8.6) {水平线：不能用};
\begin{scope}[shift={(11.85,5.3)},scale=.43]
\draw[teal!75!black,thick,domain=-2.5:4.2,samples=80,variable=\v] plot ({\v*\v/4},{\v});
\draw[->,gray] (-.7,0)--(4.7,0);
\fill (1,0) circle(5pt);\node at (2.1,-.7) {$F$};
\draw[violet,thick] (-.7,0)--(4.7,0);\fill[red!75!black] (0,0) circle(6pt);\node[above left] at (0,0) {$O$};
\end{scope}
\node[align=center] at (11.85,2.8) {$y=0$，只有顶点 $O$\\[12pt]没有第二个有限端点};
\node at (7.1,.65) {$y^2=4x$：通径既能校验 $1/p+1/p=2/p$，也能取到最短弦长 $2p$。};
\end{tikzpicture}
\end{center}

通径 $x=p/2$ 是真实焦点弦，整条长 $2p$；每半条长 $p$。
水平线虽过焦点，却只有一个有限交点，不能代入两端点公式。
角度式在 $\cos\theta=\pm1$ 时出现零分母，只表示远端趋向无穷的极限。
\end{explanationbox}
\clearpage
\begin{explanationbox}
\textbf{第五步：开口可以改变，正参数与长度定值不变。}
\begin{center}\begin{tikzpicture}[font=\small,>=stealth]
\draw[gray!45,rounded corners=3pt] (0,0) rectangle (14.2,14);
\node at (7.1,13.35) {对称或旋转只改变开口，不改变两段长度};
\node at (3.6,12.1) {向右：$y^2=2px$};
\begin{scope}[shift={(3.6,9.6)},scale=.4]
\draw[->,gray] (-4.9,0)--(4.9,0);
\draw[->,gray] (0,-4.9)--(0,4.9);
\draw[teal!75!black,thick,domain=-2.2:4.2,samples=80,variable=\v] plot ({\v*\v/4},{\v});
\draw[blue!75!black,very thick] (1, 0)--(4, 4);
\draw[orange!85!black,very thick] (1, 0)--(0.25, -1);
\fill (1, 0) circle(5pt);\node[below right] at (1, 0) {$F$};
\fill (4,4) circle(5pt);\fill (.25,-1) circle(5pt);
\end{scope}
\node at (10.5,12.1) {向左：$y^2=-2px$};
\begin{scope}[shift={(10.5,9.6)},scale=.4]
\draw[->,gray] (-4.9,0)--(4.9,0);
\draw[->,gray] (0,-4.9)--(0,4.9);
\draw[teal!75!black,thick,domain=-2.2:4.2,samples=80,variable=\v] plot ({-\v*\v/4},{\v});
\draw[blue!75!black,very thick] (-1, 0)--(-4, 4);
\draw[orange!85!black,very thick] (-1, 0)--(-0.25, -1);
\fill (-1, 0) circle(5pt);\node[below left] at (-1, 0) {$F$};
\fill (-4,4) circle(5pt);\fill (-.25,-1) circle(5pt);
\end{scope}
\node at (3.6,6.1) {向上：$x^2=2py$};
\begin{scope}[shift={(3.6,3.5999999999999996)},scale=.4]
\draw[->,gray] (-4.9,0)--(4.9,0);
\draw[->,gray] (0,-4.9)--(0,4.9);
\draw[teal!75!black,thick,domain=-2.2:4.2,samples=80,variable=\v] plot ({-\v},{\v*\v/4});
\draw[blue!75!black,very thick] (0, 1)--(-4, 4);
\draw[orange!85!black,very thick] (0, 1)--(1, 0.25);
\fill (0, 1) circle(5pt);\node[above right] at (0, 1) {$F$};
\fill (-4,4) circle(5pt);\fill (1,.25) circle(5pt);
\end{scope}
\node at (10.5,6.1) {向下：$x^2=-2py$};
\begin{scope}[shift={(10.5,3.5999999999999996)},scale=.4]
\draw[->,gray] (-4.9,0)--(4.9,0);
\draw[->,gray] (0,-4.9)--(0,4.9);
\draw[teal!75!black,thick,domain=-2.2:4.2,samples=80,variable=\v] plot ({\v},{-\v*\v/4});
\draw[blue!75!black,very thick] (0, -1)--(4, -4);
\draw[orange!85!black,very thick] (0, -1)--(-1, -0.25);
\fill (0, -1) circle(5pt);\node[below left] at (0, -1) {$F$};
\fill (4,-4) circle(5pt);\fill (-1,-.25) circle(5pt);
\end{scope}
\node at (7.1,.65) {四图均取正参数 $p=2$；蓝段长 $5$、橙段长 $5/4$，倒数和均为 $1$。};
\end{tikzpicture}
\end{center}

四式统一规定 $p>0$。若写成 $y^2=4ax$（$a>0$），则 $p=2a$，倒数和为 $1/a$。
\textbf{识别顺序：}先确认两个不同有限交点且弦过焦点，再用倒数和；求弦长时配合 $AB=r+s$。
\end{explanationbox}
